\documentclass[11pt]{article}
\usepackage[margin=1in]{geometry}
\usepackage{amsmath,amssymb,amsthm}
\usepackage[colorlinks=true,urlcolor=blue]{hyperref}
\newtheorem{theorem}{Theorem}
\newtheorem{lemma}{Lemma}
\newtheorem{corollary}{Corollary}
\newcommand{\Fp}{\mathbb F_p}
\newcommand{\R}{\mathbb R}
\newcommand{\one}{\mathbf1}
\newcommand{\Prob}{\mathbb P}
\newcommand{\E}{\mathbb E}
\newcommand{\nullity}{\operatorname{nullity}}
\newcommand{\rank}{\operatorname{rank}}
\title{Configuration-model strong mixing and the odd-block RSBM extension}
\author{Partitioning frontier session: proof synthesis}
\date{November 15, 2026}
\begin{document}
\maketitle
\begin{abstract}
We combine three fully proved companion estimates to obtain the
summed maximum-atom mixing law over every fixed odd prime field for the
undirected regular configuration model, along every admissible size.
The companions give a direct nearconstant stub-exposure bound, an
arbitrary-target entropy estimate, and a relative central local limit
theorem with its lattice factor. This is an explicitly supplied
extension of the counting method of Huang and the mixing conclusion
of M\'esz\'aros; it is not a claim that their stated theorems already
cover arbitrary targets on odd vertex sets. The resulting two-block
first moment proves simplicity of the planted RSBM eigenvalue for
every fixed internal degree at least three, including odd block size
when the internal degree is even. Historical novelty and effective
failure rates are not established.
\end{abstract}

\section{The configuration mixing statement}
Let $d\ge3$ be fixed. At each of $n$ labelled vertices put $d$ labelled
half-edges, pair the $dn$ half-edges uniformly, and let $C_n$ count
paired half-edges between vertices. Thus a loop contributes two to
the diagonal and $C_n\one=d\one$. The only size restriction is that
$dn$ is even. Fix an odd prime $p$, independently of $n$, and work
over $\Fp$.
For $q\in\Fp^n$ write $s(q)=\sum_iq_i$ and define
\[
 Q_q(r)=\begin{cases}
  \boldsymbol1_{\{r=da\one\}},&q=a\one,\\
  p^{-n+1}\boldsymbol1_{\{s(r)=ds(q)\}},&q\text{ nonconstant}.
 \end{cases}
\]
The hyperplane in the second line has exactly $p^{n-1}$ elements,
including when $p\mid n$ or $p\mid d$.

\begin{theorem}[All-admissible-size strong mixing]\label{thm:mix}
For every fixed $d\ge3$ and fixed odd prime $p$,
\[
 e_n:=\sum_{q\in\Fp^n}\sup_{r\in\Fp^n}
       \bigl|\Prob(C_nq=r)-Q_q(r)\bigr|\longrightarrow0
 \qquad(dn\text{ even},\ n\to\infty).
\]
No prime growing with $n$ is asserted. No evenness assumption on $n$
is used when $d$ is even.
\end{theorem}

\section{Precisely which companion results are used}
All three companion notes contain their proofs, rather than citing
a zero-target invertibility theorem as a shifted or all-target result.
For a vector $q$, set
\[
 n_a(q)=|\{i:q_i=a\}|,\qquad \nu_a(q)=n_a(q)/n,
 \qquad u=(1/p,\ldots,1/p).
\]

\begin{lemma}[Nearconstant inputs]\label{lem:near}
There is $\eta=\eta(d,p)>0$ such that
\[
 \sum_{\substack{q\text{ nonconstant}\\
             \max_a\nu_a(q)\ge1-\eta}}
       \sup_r\Prob(C_nq=r)=O_{d,p}(n^{1-d/2})=o(1).
\]
\end{lemma}
The complete elementary proof is in the companion note:
\begin{quote}\small\ttfamily
parity-reduction/\\
nearconstant-configuration-atoms.tex
\end{quote}
For $q=a\one+z$, $|\operatorname{supp}z|=k$, it exposes the partners
of stubs in that support. For any prescribed target, at least $dk/2$
exposures must hit a prescribed or previously hit vertex. It proves
\[
 \sup_r\Prob(C_nz=r)
 \le2^{dk+1}(K_d k/n)^{dk/2},\qquad K_d=2(2d+1),
\]
for sufficiently small $k/n$. Summing
$p\binom nk(p-1)^k$ inputs proves Lemma~\ref{lem:near}.

\begin{lemma}[Noncentral nonconstant inputs]\label{lem:middle}
For every fixed $\eta>0$ and $A>0$, there is a finite
$B=B(d,p,\eta,A)$ such that
\[
 \sum_{\substack{q:\max_a\nu_a(q)\le1-\eta\\
                \|\nu(q)-u\|_2^2>B\log(n+1)/n}}
       \sup_r\Prob(C_nq=r)=O_{d,p,\eta,A}(n^{-A}).
\]
\end{lemma}
The full proof is the all-target noncentral-tail theorem in the companion note:
\begin{quote}\small\ttfamily
huang-primary/\\
configuration-arbitrary-target-tails.tex
\end{quote}
It uses an exact prescribed-output extension of Huang's
Proposition~2.2~\cite{Huang}. For clarity, its crucial inequality is
proved there for every target:
\[
 F_t(\pi):=\sum_{a_1+\cdots+a_d=t}
          \prod_j\pi(a_j)^{(d-1)/d}\le1.
\]
Equality is possible only at the uniform distribution or a point mass.
For symmetric edge-color probability matrix $m$ with marginal $\nu$,
the summed type rate is bounded by
\[
 \Psi(m)=-\frac{d-2}{2}D(m\|\nu\otimes\nu)
         +\sum_a\nu_a\log\max_tF_t(m_a/\nu_a).
\]
The note proves the uniform expansion
\[
 \log F_t(u+z)=-\frac{(d-1)p}{2d}\|z\|_2^2+O_{d,p}(\|z\|_2^3)
 \quad\left(\sum_a z_a=0\right).
\]
Consequently $\Psi$ has a strictly negative quadratic bound at the
uniform matrix; compactness gives the same bound away from the
point-mass neighborhoods. The polynomial factors and count of type
matrices are included. This proves Lemma~\ref{lem:middle}, uniformly
over targets. It is not a zero-output-only large-deviation estimate.

\begin{lemma}[Balanced arbitrary-target local limit]\label{lem:central}
For every fixed $K>0$, uniformly over $q$ with
\[
 \max_a|n_a(q)-n/p|\le K\sqrt{n\log n},
\]
and over every $r$,
\[
 \Prob(C_nq=r)=
 \begin{cases}
 (1+o(1))p^{-n+1},&s(r)=ds(q),\\
 0,&\text{otherwise}.
 \end{cases}
\]
\end{lemma}
The complete proof is Theorem~1 of the companion note:
\begin{quote}\small\ttfamily
meszaros-extension/\\
balanced-all-target-configuration.tex
\end{quote}
It counts ordered partner colors at each vertex, conditioned on that
vertex's prescribed output. All such $d$-tuple count distributions
have the same mean and covariance, because any two positions are
independent uniform labels for $d\ge3$. Exponential tilting supplies
a \emph{relative} moderate-deviation local limit estimate, uniform
over the entire sequence of prescribed outputs.
Each row contributes a lattice factor $p$. Symmetry leaves $p-1$
independent weighted-row congruences, whose cosets have equal central
pairing mass. The explicit bounded translations
\[
 Z_j=2(e_j-e_0)(e_j-e_0)^T,
 \qquad \Phi(Z_j)=2j(e_j-e_0),\qquad j\ne0,
\]
prove the index $p^{p-1}$, including the even-diagonal restriction.
The net factor is $p^p/p^{p-1}=p$, producing $p^{-n+1}$.
The proof handles central-type tails and does not assume parity
independence or a Gaussian leading constant without computing it.

\section{Assembly of the full strong-mixing law}
\begin{proof}[Proof of Theorem~\ref{thm:mix}]
Constant inputs have zero error, exactly. Shrink the $\eta$ in
Lemma~\ref{lem:near} if necessary so that
\[
 0<\eta<\tfrac12(1-1/p),\qquad
 h(\eta)+\eta\log(p-1)<\log p,
\]
where $h(t)=-t\log t-(1-t)\log(1-t)$. This is possible since the
left side of the second inequality tends to zero with $\eta$.

Partition the remaining inputs into a nearconstant class
$\mathcal N_n$ with $\max_a\nu_a\ge1-\eta$, a central class
$\mathcal C_n$ with
$\|\nu-u\|_2^2\le B\log(n+1)/n$, and the remainder
$\mathcal I_n$. For sufficiently large $n$, the first two classes are
disjoint. Choose $B$ from Lemma~\ref{lem:middle}, with $A=2$.

For any nonconstant input, $\sup_r Q_q(r)=p^{-n+1}$. Thus, by
Lemma~\ref{lem:near}, the total error on $\mathcal N_n$ is at most
\[
 o(1)+p^{-n+1}|\mathcal N_n|=o(1).
\]
Indeed $|\mathcal N_n|\le p\sum_{k\le\eta n}\binom nk(p-1)^k$
is exponentially smaller than $p^n$ by the chosen entropy inequality.

Lemma~\ref{lem:middle} bounds the actual probabilities summed over
$\mathcal I_n$ by $O(n^{-2})$. For the reference probabilities,
\[
 p^{-n+1}|\mathcal I_n|
 \le p\Prob\!\left(\|\widehat\nu-u\|_2^2>
                       B\log(n+1)/n\right)=o(1),
\]
where $\widehat\nu$ is the empirical profile of $n$ independent
uniform field elements. The last limit follows, for example, from
Hoeffding's inequality and the union bound over the fixed $p$ labels.
Hence the total error on $\mathcal I_n$ tends to zero.

For $q\in\mathcal C_n$, every coordinate count differs from $n/p$
by at most $\sqrt{Bn\log(n+1)}$. Lemma~\ref{lem:central}, with a
fixed $K>\sqrt B$, gives
\[
 \sup_{q\in\mathcal C_n,r}
       |\Prob(C_nq=r)-Q_q(r)|=o(p^{-n+1}).
\]
There are at most $p^n$ such inputs, so their total error is $o(1)$
as well. All constants keep $p,d$ fixed. This proves the theorem
along every sequence with $dn$ even.
\end{proof}

\section{Two planted blocks, including odd sizes}
Fix integers $d_1>d_2>0$ with $d_1\ge3$, and let
$\delta=d_1-d_2$. For a simple $d_2$-regular bipartite graph on two
blocks of size $N$, let $B$ be its $N$-by-$N$ zero-one biadjacency
matrix. Initially let $C,D$ be independent degree-$d_1$
configuration-model adjacency matrices, and put
\[
 A_B=\begin{pmatrix}C&B\\B^T&D\end{pmatrix},\qquad
 \sigma=\binom{\one_N}{-\one_N}.
\]
Always $A_B\sigma=\delta\sigma$.

\begin{theorem}[All-admissible-size planted simplicity]\label{thm:rsbm}
Let $C,D$ instead be independent uniform \emph{simple}
$d_1$-regular graphs on their respective labelled blocks. Then
\[
 \lim_{\substack{N\to\infty\\d_1N\text{ even}}}
 \sup_{B}\Prob\bigl(\nullity_{\R}(A_B-\delta I)\ge2\bigr)=0,
\]
where the supremum is over all simple $d_2$-regular bipartite
biadjacency matrices. Thus the result also holds for an independent
uniform simple cross graph. In particular, even $d_1\ge4$ admits
\emph{both} odd and even block sizes. No outlier spectral gap is
assumed.
\end{theorem}

\begin{proof}
Fix an odd prime $p\nmid d_2$, and apply
Theorem~\ref{thm:mix} to degree $d_1$, at block size $N$.
For one block define
\[
 e_x=\sup_r|\Prob(Cx=r)-Q_x(r)|,\qquad
 e_N=\sum_xe_x=o(1).
\]
For every fixed $b\in\Fp^N$,
\begin{equation}\label{eq:affine}
 \sum_xQ_x(\delta x+b)\le2,
 \qquad \sum_x\Prob(Cx=\delta x+b)\le2+e_N.
\end{equation}
For nonconstant $x$, the reference sum condition becomes
$d_2s(x)=s(b)$, allowing at most $p^{N-1}$ inputs of mass
$p^{-N+1}$. For constant $x=a\one$, it requires
$b=d_2a\one$, allowing at most one constant. This proves
\eqref{eq:affine} uniformly in $b$.

Independence of the two internal configurations gives
\[
 \E|\ker_{\Fp}(A_B-\delta I)|
 =\sum_{x,y}\Prob(Cx=\delta x-By)
                \Prob(Dy=\delta y-B^Tx).
\]
Replace the first probability by $Q_x$, then the second by $Q_y$.
By \eqref{eq:affine} the total error is at most
$(2+e_N)e_N+2e_N=4e_N+e_N^2$, uniformly in $B$.
The remaining reference sum is bounded as follows.

If $x,y$ are both nonconstant, the two hyperplane conditions reduce
to $s(x)+s(y)=0$ because $B\one=B^T\one=d_2\one$ and
$p\nmid d_2$. At most $p^{2N-1}$ pairs, each of mass
$p^{-2N+2}$, contribute, giving at most $p$.
If both are constant, writing $x=a\one,y=b\one$ gives exactly
$a+b=0$, hence contribution $p$.
If only $x=a\one$ is constant, the first equation requires
$By=-d_2a\one$. For each $a$ there are at most
$p^{N-r_B}$ solutions, where $r_B=\rank_{\Fp}B$.
The second factor has mass at most $p^{-N+1}$, giving at most
$p^{2-r_B}$; the opposite orientation gives the same bound.
Consequently
\[
 \E|\ker_{\Fp}(A_B-\delta I)|
 \le2p+2p^{2-r_B}+4e_N+e_N^2.
\]
Greedily select rows of $B$ with disjoint supports. Each selected
row excludes at most $1+d_2(d_2-1)$ rows. The selected nonzero rows
are independent over every field, so
\[
 r_B\ge N/[1+d_2(d_2-1)].
\]
Thus the limsup of the expected kernel size is at most $2p$,
uniformly in $B$, along all admissible $N$.

If the real nullity of the integer matrix is at least two, all its
integer minors of order $2N-1$ vanish, so its mod-$p$ nullity is also
at least two. Therefore
\[
 \limsup_{\substack{N\to\infty\\d_1N\text{ even}}}
 \sup_B\Prob\bigl(\nullity_{\R}(A_B-\delta I)\ge2\bigr)
 \le2/p
\]
in the configuration model. Taking arbitrarily large \emph{fixed}
odd primes not dividing $d_2$ makes this limsup zero.

For fixed $d_1$, simplicity of a configuration has probability
bounded below away from zero along all admissible sizes, and
conditioning it on simplicity gives a uniform labelled simple
$d_1$-regular graph. The two independent configurations can both be
conditioned at bounded cost. The bound already holds uniformly in
$B$, so conditioning does not change that uniformity or the independence
of a subsequently sampled cross graph. This proves the theorem.
\end{proof}

\begin{corollary}[Exact algebraic recovery]
For the parameter range of Theorem~\ref{thm:rsbm}, exact rational
elimination of $A-(d_1-d_2)I$ returns the planted signs, up to global
exchange, with probability tending to one, in polynomial bit complexity.
\end{corollary}
\begin{proof}
On the one-dimensional-kernel event, any nonzero rational kernel
vector is a multiple of $\sigma$. Divide by its first coordinate.
The integer rows have bounded norm for fixed degrees; Hadamard's
inequality bounds minors by numbers with $O(N)$ bits. Standard
fraction-free elimination therefore has polynomial bit complexity.
\end{proof}

\paragraph{Remaining scope and provenance.}
The $(d_1,d_2)=(2,1)$ case is not covered: all three mixing estimates
use $d_1\ge3$. The proof gives qualitative high probability, not a
uniform quantitative rate in a growing prime. The familiar planted
eigenvector identity and exact elimination are not claimed new.
The new proof burden addressed here is the all-target configuration
mixing law across the parity restriction, with its actual leading
constant and charged target uniformity. The imported primary statements
of Huang and M\'esz\'aros are narrower than that conclusion; the
companion proofs explicitly supply the extensions. Historical firstness
remains unverified.

\begin{thebibliography}{9}
\bibitem{Huang} J.~Huang, Invertibility of adjacency matrices for random
$d$-regular graphs, arXiv:1807.06465v2 (2018), Proposition~2.2 and
Section~4. \url{https://arxiv.org/pdf/1807.06465v2}.
\bibitem{Meszaros} A.~M\'esz\'aros, The distribution of sandpile groups
of random regular graphs, arXiv:1806.03736v4 (2020), Theorem~1.6.
\url{https://arxiv.org/pdf/1806.03736v4}.
\bibitem{Bollobas} B.~Bollob\'as, A probabilistic proof of an asymptotic
formula for the number of labelled regular graphs, European
J.~Combin.~1 (1980), 311--316.
\url{https://doi.org/10.1016/S0195-6698(80)80030-8}.
\end{thebibliography}
\end{document}
